\scanpage{41}

\[
\mathrm{Hom}^{d}(\tau_\alpha M_{S_k},\underset{\substack{\cap\\\tau_{\alpha+d}N_{S_k}}}{\underline{\tau_\beta N_{S_k}}})
\hookrightarrow\mathrm{Hom}^{d}(\tau_\alpha M_{S_k},\tau_{\alpha+d}N_{S_k})
\]

% The oblique leader starts at the left Hom in the preceding display.
\inserted{(sous-mod. des hom de degré $d$ de $\tau_\alpha M_{S_k}$ \unsure{ds|dans} $\tau_{\alpha+d}N_{S_k}$ qui veulent bien appliquer $\tau_\alpha M_{S_k}$ dans $\tau_\beta N_{S_k}\subset\tau_{\alpha+d}N_{S_k}$)}

\[
\overset{\rotatebox{90}{$\simeq$}}{\mathrm{Hom}_k(M,N)}
\]

Il faut donc prendre les $u:M\longrightarrow N$ tels que l'hom correspondant $\tau_\alpha M_{S_k}\mathord{-}\tau_{\alpha+d}N_{S_k}$

\[
mT^{(j)}\longmapsto C_{d,j}\,u(m)T^{(d+j)}
\]

\deleted{applique $M_S$} a son image dans $\tau_\beta N_{S_k}$, i.e. tels que

\[
\left\{\begin{aligned}
C_{d,\alpha}u&=0\\
C_{d,\alpha+1}u&=0\\
&\cdots\\
C_{d,\beta-d}u&=0
\end{aligned}\right.
\]

(ce qui s'écrit

\[
C_{d,j}\,u=0\quad\text{si}\quad\alpha\leq j<\underset{\deleted{\unsure{\beta-d+1}}}{\beta-d}\qquad\bigg)
\]

c) \uline{$d<0$}. Notons que

\[
N^{*}_{\deleted{\otimes}S_k}=\varphi_{0\beta}\tau_\beta N_{S_k}=\tau_0\underset{\substack{\cap\\\varphi_\beta N_{S_k}}}{\varphi_\beta\tau_\beta N_{S_k}}
\]

donc \inserted{si $M^{*}\in\mathrm{Ob}\,\mathcal{M}^{\unsure{*|\prime}}$}

\[
\begin{aligned}
\mathrm{Hom}^{d}(M^{*},N_{S_k})&\subset\mathrm{Hom}^{d}(M^{*},\varphi_\beta\tau_\beta N_{S_k})\\
&\hspace{9mm}\shortparallel\\
&\mathrm{Hom}(M^{\unsure{\prime|*}}[-d],\varphi_\beta N_{S_k})\\
&\hspace{9mm}\shortparallel\\
&\mathrm{Hom}(\tau_\beta M^{\prime}[-d],N_{S_k}
\end{aligned}
\]

% archive sheet number: 40
